% =====================================================================
%  09 附录
%  正文外的支撑内容:核心代码、补充数据、长表格等。
%  由 main.tex 的 appendices 环境包裹,放在正文最后。
%  代码用 lstlisting 环境(模板已加载 listings 宏包)。
% =====================================================================

\section{核心代码}

\subsection{问题一:半经验关联式标定与 GPR 交叉验证}

\begin{lstlisting}[language=python, basicstyle=\small\ttfamily, numbers=left, numberstyle=\tiny]
# --- 压降幂律关联式 ln(dP*) = lnC + a*ln(beta) + b*ln(1-alpha) + c*ln(n) ---
def fit_dp_corr(df):
    y = np.log(df["dP"].to_numpy())
    X = np.column_stack([
        np.ones(len(df)),
        np.log(df["beta"].to_numpy()),
        np.log(1 - df["alpha"].to_numpy()),
        np.log(df["n"].to_numpy()),
    ])
    coef, *_ = np.linalg.lstsq(X, y, rcond=None)
    lnC, a, b, c = coef
    C = np.exp(lnC)
    dp_pred = C * df["beta"]**a * (1 - df["alpha"])**b * df["n"]**c
    r2 = r2_score(df["dP"], dp_pred)
    mape = mean_absolute_percentage_error(df["dP"], dp_pred)
    return {"C": C, "a": a, "b": b, "c": c, "r2": r2, "mape": mape, "pred": dp_pred}

# --- 热阻二次结构关联式 Rth* = R0 + c1*(alpha-alpha*)^2 + c2*beta + c3*n + c4*alpha*n ---
def rth_model(X, R0, c1, c2, c3, c4, alpha_star):
    alpha, beta, n = X
    return R0 + c1*(alpha-alpha_star)**2 + c2*beta + c3*n + c4*alpha*n

def fit_rth_corr(df):
    X = (df["alpha"].to_numpy(), df["beta"].to_numpy(), df["n"].to_numpy())
    y = df["Rth"].to_numpy()
    p0 = [0.75, 0.1, 0.0, 0.0, 0.0, 0.2]
    popt, _ = curve_fit(rth_model, X, y, p0=p0, maxfev=20000)
    R0, c1, c2, c3, c4, alpha_star = popt
    rth_pred = rth_model(X, *popt)
    r2 = r2_score(y, rth_pred)
    mape = mean_absolute_percentage_error(y, rth_pred)
    return {"R0": R0, "c1": c1, "c2": c2, "c3": c3, "c4": c4,
            "alpha_star": alpha_star, "r2": r2, "mape": mape, "pred": rth_pred}

# --- alpha* GPR交叉验证:在 GPR 上扫描 alpha 求 Rth 均值 argmin ---
def gpr_alpha_star(gprs, scaler, beta_grid, n_grid, n_points=41):
    alphas = np.linspace(0.10, 0.30, n_points)
    avg = []
    for a in alphas:
        vals = [gprs["Rth"].predict(scaler.transform([[a, b, n]]))[0]
                for b in beta_grid for n in n_grid]
        avg.append(np.mean(vals))
    i = np.argmin(avg)
    return alphas[i], avg[i]
\end{lstlisting}

\subsection{问题二:GPR 代理模型构建(Constant+RBF+White 复合核)}

\begin{lstlisting}[language=python, basicstyle=\small\ttfamily, numbers=left, numberstyle=\tiny]
# --- 数据划分与标准化 ---
X_train, X_test, y_train, y_test = train_test_split(
    data[["alpha","beta","n"]], data[["Rth","dP","dT"]],
    test_size=0.2, random_state=676767)
scaler = StandardScaler()
X_train_std = scaler.fit_transform(X_train)
X_test_std  = scaler.transform(X_test)

# --- 复合核: Constant + RBF(各向异性) + White ---
kernel = (ConstantKernel(1.0, (1e-10, 1e10))
          * RBF([1.0, 1.0, 1.0], (1e-10, 1e10))
          + WhiteKernel(1.0, (1e-10, 1e10)))

# --- 对每个输出维度独立训练 GPR (MISO 架构) ---
result = {}
for y_col in ["Rth", "dP", "dT"]:
    gpr = GaussianProcessRegressor(
        kernel=kernel, alpha=1e-8,
        normalize_y=True, n_restarts_optimizer=10)
    gpr.fit(X_train_std, y_train[y_col])
    mu, sigma = gpr.predict(X_test_std, return_std=True)

    # 提取 MLE 最优超参数 theta* = [sigma_f, l1, l2, l3, sigma_n]
    k_cr, k_w = gpr.kernel_.k1, gpr.kernel_.k2
    sigma_f = np.sqrt(k_cr.k1.constant_value)
    length_scales = k_cr.k2.length_scale
    sigma_n = np.sqrt(k_w.noise_level)

    result[y_col] = {
        "gpr": gpr, "y_true": y_test[y_col].to_numpy(),
        "y_pred": mu, "y_std": sigma,
        "r2"  : r2_score(y_test[y_col], mu),
        "rmse": np.sqrt(mean_squared_error(y_test[y_col], mu)),
        "mape": mean_absolute_percentage_error(y_test[y_col], mu),
        "theta_star": {"sigma_f": sigma_f,
                       "length_scales": length_scales,
                       "sigma_n": sigma_n},
    }

result["_scaler"] = scaler
joblib.dump(result, DATA_DIR / "gpr.pkl")
\end{lstlisting}

\subsection{问题三:NSGA-II + L-BFGS-B 混合多目标优化与理想点法决策}

\begin{lstlisting}[language=python, basicstyle=\small\ttfamily, numbers=left, numberstyle=\tiny]
# --- 混合变量 NSGA-II 问题定义(mode 变量处理针肋排数耦合约束) ---
class SinkOptProblem(Problem):
    def __init__(self):
        super().__init__(
            vars={
                "mode": Choice(options=[0, 1]),        # 0=无针肋, 1=有针肋
                "x1"  : Real(bounds=(0.10, 0.30)),     # 针肋宽度比
                "x2"  : Choice(options=[3.0, 3.5, 4.0, 4.5]),  # 歧管深高比(四档)
                "x3"  : Integer(bounds=(2, 10)),       # 针肋排数
            },
            n_obj=3,
        )
    def _evaluate(self, X, out, *args, **kwargs):
        arr = np.array([to_physical(x) for x in X])
        out["F"] = predict(arr)  # GPR 预测 [Rth, dP, dT]

# --- 多种子拼接 Pareto 前沿 ---
def run_nsga2_multiseed(pop_size=200, n_gen=300, seeds=range(10)):
    X_all, F_all = [], []
    for seed in seeds:
        problem = SinkOptProblem()
        algorithm = NSGA2(pop_size=pop_size,
            sampling=MixedVariableSampling(),
            mating=MixedVariableMating(
                eliminate_duplicates=MixedVariableDuplicateElimination()),
            eliminate_duplicates=MixedVariableDuplicateElimination())
        result = pymoo_minimize(problem, algorithm,
            termination=("n_gen", n_gen), seed=seed, verbose=False)
        X_all.append(np.array([to_physical(x) for x in result.X]))
        F_all.append(result.F)
    return np.vstack(X_all), np.vstack(F_all)

# --- 理想点法:固定 (x2,x3) → L-BFGS-B 精修 x1,穷举全局最优 ---
def refine_at_fixed_x2x3(x2, x3, x1_init, f_min, f_max, w):
    if x3 == 0:
        return 0.0, ideal_dist(np.array([0.0, x2, x3]), f_min, f_max, w)
    def obj(x1):
        F = predict(np.array([[x1[0], x2, x3]]))[0]
        F_norm = (F - f_min)/(f_max - f_min)
        return np.sqrt(np.sum(w * F_norm**2))
    result = lbfgsb(obj, x0=[x1_init], bounds=[(0.10, 0.30)],
                    method="L-BFGS-B")
    return result.x[0], result.fun

def global_best(w, f_min, f_max, x1_inits=(0.10,0.15,0.20,0.25,0.30)):
    best = None
    for x2, x3 in ALL_COMBOS:          # 40 种合法 (x2,x3) 组合
        inits = (0.0,) if x3 == 0 else x1_inits
        for x1_init in inits:
            x1_r, D_r = refine_at_fixed_x2x3(x2, x3, x1_init,
                                              f_min, f_max, w=w)
            if best is None or D_r < best[3]:
                best = (x1_r, x2, x3, D_r)
    return best  # (x1*, x2*, x3*, D*)
\end{lstlisting}

\subsection{问题四:最小最大后悔值鲁棒优化}

\begin{lstlisting}[language=python, basicstyle=\small\ttfamily, numbers=left, numberstyle=\tiny]
# --- 权重单纯形网格采样 ---
def simplex_grid(n=6):
    pts = []
    for i in range(n+1):
        for j in range(n+1-i):
            pts.append([i/n, j/n, (n-i-j)/n])
    return np.array(pts)

# --- Minimax Regret 鲁棒解搜索 ---
def minimax_regret(candidates, weight_grid, f_min, f_max, D_star_grid):
    F_cand = predict(candidates)               # (N, 3)
    F_norm = (F_cand - f_min)/(f_max - f_min)  # Min-Max 归一化

    max_regret = np.full(len(candidates), -np.inf)
    for w, D_opt in zip(weight_grid, D_star_grid):
        D_w = np.sqrt(np.sum(w * F_norm**2, axis=1))
        regret = D_w - D_opt
        max_regret = np.maximum(max_regret, regret)

    robust_idx = np.argmin(max_regret)
    return robust_idx, max_regret

# --- 候选方案的跨场景稳定性评估(变异系数) ---
def cv_of_candidate(X, weight_grid, f_min, f_max):
    F = predict(np.array([X]))[0]
    F_norm = (F - f_min)/(f_max - f_min)
    D_w = np.sqrt(np.sum(weight_grid * F_norm**2, axis=1))
    return D_w.std() / D_w.mean()
\end{lstlisting}

\subsection{问题五:Sobol 全局敏感性 + LHS 蒙特卡洛不确定性传播}

\begin{lstlisting}[language=python, basicstyle=\small\ttfamily, numbers=left, numberstyle=\tiny]
# --- Sobol 全局敏感性指数(一阶 S_i + 总效应 S_Ti, N=1024) ---
SOBOL_PROBLEM = {
    "num_vars": 3, "names": ["x1","x2","x3"],
    "bounds": [[0.10, 0.30], [3.00, 4.50], [2, 10]]}

def sobol_indices(N=1024):
    param_values = sobol_sample.sample(SOBOL_PROBLEM, N)
    result = {}
    for col in ["Rth", "dP", "dT"]:
        X_std = scaler.transform(param_values)
        Y = gprs[col].predict(X_std)
        Si = sobol_analyze.analyze(SOBOL_PROBLEM, Y, print_to_console=False)
        result[col] = {"S1": Si["S1"], "ST": Si["ST"]}
    return result

# --- 局部归一化敏感性系数(中心差分,标称点 X0 处) ---
def lsa_structural(X0, h_rel=1e-3):
    F0 = predict(X0)[0]
    S = np.zeros((3, 3))
    for i in range(3):
        h = max(abs(X0[i]) * h_rel, 1e-4)
        Xp, Xm = X0.copy(), X0.copy()
        Xp[i] += h; Xm[i] -= h
        dF = (predict(Xp)[0] - predict(Xm)[0])/(2*h)
        S[i] = dF * X0[i] / F0
    return S  # S[i,j] = dGj/dxi * xi0/yj0

# --- LHS 蒙特卡洛不确定性传播(结构参数层, N_MC=8000, eps=5%) ---
def lhs_montecarlo(X0, eps=0.05, N=8000, seed=676767):
    sigma = eps * np.abs(X0) / 3   # 3-sigma 落在 +-5% 内
    sampler = qmc.LatinHypercube(d=3, seed=seed)
    U = sampler.random(N)
    Z = norm.ppf(U)                 # 标准正态分位数 → N(0,1)
    samples = X0 + Z * sigma        # 扰动样本(N, 3)
    F = predict(samples)            # GPR 传播(N, 3)
    F0 = predict(X0)[0]
    cv = F.std(axis=0, ddof=1) / F.mean(axis=0)
    y_limit = 1.10 * F0             # 越界阈值:+10%
    p_fail = (F > y_limit).mean(axis=0)
    return {"F0": F0, "cv": cv, "p_fail": p_fail, "F_samples": F}

# --- 工况参数层:基于无量纲标度关系的机理线性化近似 ---
# mdot→Rth: S≈-1/3 (Nu~Re^1/3), mdot→dP: S≈-1 (f~1/Re)
# Tin: 一阶抵消 S≈0; qv: 分子分母同步缩放 S≈0
OP_SENSITIVITY = {
    "mdot": {"Rth": 0.0, "dP": -1.0, "dT": 0.0},
    "Tin":  {"Rth": 0.0, "dP":  0.0, "dT": 0.0},
    "qv":   {"Rth": 0.0, "dP":  0.0, "dT": 0.0},
}
\end{lstlisting}
